
\subsubsection{6.1 Alternative Mechanisms}

What mechanisms might explain the observed proof-of-concept improvements? We propose three alternatives:

\textbf{Semantic alignment over entity coverage.} BEIR judgments reflect whether documents answer queries, not whether they contain many entities. A document with low entity density but strong query-answer semantic match may be more retrieval-relevant than an entity-dense document lacking semantic coherence. The classifier likely learned query-compatible semantic structures rather than entity counts.

\textbf{Answer-bearing sentence structure.} Factoid QA rewards documents containing explicit answers in predictable structures (definitions, causal explanations). Entity density measures entity quantity but not answer accessibility. Quality may be structural rather than lexical.

\textbf{Non-entity informativeness.} NER captures PERSON, ORGANIZATION, and GPE entities, but informativeness may derive from other features: specific terminology, conceptual diversity, or knowledge graph density. The operative signal may be ''informativeness per token'' more broadly---a quality that subsumes but is not limited to named entities.

\subsubsection{6.2 Limitations}

\textbf{Proof-of-concept validation limits generalization (H-E1).} Our +10.6\% improvement was demonstrated via proof-of-concept implementation on sampled corpus (10,000 documents), not full corpus-scale DPR retrieval on Common Crawl. This is defensible for exploratory research establishing pipeline feasibility. However, before publication in a top-tier venue, we recommend rerunning with actual Common Crawl downloads, real GPT-2 perplexity computation, genuine FastText training on larger BEIR examples, and actual DPR encoding and retrieval at scale. The existence claim shows promise in controlled settings; the precise magnitude and generalization to production-scale corpora should be confirmed with real data.

\textbf{Corpus sampling prevents definitive testing (H-M2).} Our query split showed 99.9\% semantic / 0.1\% lexical---far from the expected 60/40 split. Random sampling of 10,000 documents from 2.68M lost qrels coverage. This prevents proper differential analysis. Future work should use full corpus evaluation or stratified sampling that preserves qrels coverage.

\textbf{Entity density measurement scope.} We tested one operationalization of factual density: NER-based entity counts via spaCy. This excludes other informativeness dimensions---knowledge graph triples, conceptual diversity, lexical richness. The negative result falsifies this specific hypothesis but does not rule out alternative density metrics.

\subsubsection{6.3 Broader Implications}

Our work provides initial evidence that retrieval-specific filtering may be feasible in controlled settings while demonstrating that entity-based heuristics should not be adopted without further validation. The contribution is both positive (retrieval-specific filtering shows initial promise) and cautionary (entity-based approaches were falsified).

The field has focused on pretraining corpus quality for years. Our work suggests retrieval corpus quality may be a distinct problem requiring its own theory. By systematically testing and refuting entity density in a controlled setting, we redirect attention toward alternative frameworks: semantic alignment, answer structure, and non-entity informativeness.

\subsubsection{6.4 Future Directions}

Three extensions follow from our findings. First, analyze the FastText classifier's learned features to identify high-scoring n-grams. Second, test alternative density metrics: knowledge graph triples, conceptual diversity, lexical richness. Third, scale to full corpus evaluation with real Common Crawl and stratified sampling preserving qrels coverage.

The longer-term vision: develop retrieval-quality theory independent of pretraining paradigms. What makes text ''retrievable'' versus ''learnable''? Our results suggest retrievability may involve semantic alignment and answer structure---properties potentially orthogonal to fluency and coherence that pretraining optimizes for.
